\section*{NeurIPS Paper Checklist}

\begin{enumerate}

\item {\bf Claims}
    \item[] Question: Do the main claims made in the abstract and introduction accurately reflect the paper's contributions and scope?
    \item[] Answer: \answerYes{}.
    \item[] Justification: The abstract and introduction state that the paper is a pre-registered diagnostic of frozen traffic STGNNs, not a new model or dataset. The scope is limited to the evaluated shifts, source cities, and masking assumptions.

\item {\bf Limitations}
    \item[] Question: Does the paper discuss the limitations of the work performed by the authors?
    \item[] Answer: \answerYes{}.
    \item[] Justification: Section~\ref{sec:discussion} states the limits of the model set, single-source training, silent zero-masking assumption, and unresolved MegaCRN temporal mechanism; Appendix~\ref{sec:appendix_reproducibility} summarizes reproducibility and asset scope.

\item {\bf Theory assumptions and proofs}
    \item[] Question: For each theoretical result, does the paper provide the full set of assumptions and a complete (and correct) proof?
    \item[] Answer: \answerYes{}.
    \item[] Justification: The only theoretical component is a first-order exposure bound. Section~\ref{sec:mechanism} states its local interpretation, and Appendix~\ref{sec:appendix_theory} gives the assumptions, derivation, and iid-dropout caveat.

\item {\bf Experimental result reproducibility}
    \item[] Question: Does the paper fully disclose all the information needed to reproduce the main experimental results of the paper to the extent that it affects the main claims and/or conclusions of the paper (regardless of whether the code and data are provided or not)?
    \item[] Answer: \answerYes{}.
    \item[] Justification: Sections~\ref{sec:setup}--\ref{sec:mechanism} describe the datasets, splits, frozen-checkpoint protocol, shifts, metrics, and interventions. Appendix~\ref{sec:appendix_reproducibility} identifies the canonical result files, registry, configs, and figure-generation commands.

\item {\bf Open access to data and code}
    \item[] Question: Does the paper provide open access to the data and code, with sufficient instructions to faithfully reproduce the main experimental results, as described in supplemental material?
    \item[] Answer: \answerYes{}.
    \item[] Justification: The submission includes evaluation scripts, configs, result tables, and the recommendation artifact. The raw traffic datasets are public third-party datasets and are cited rather than redistributed.

\item {\bf Experimental setting/details}
    \item[] Question: Does the paper specify all the training and test details (e.g., data splits, hyperparameters, how they were chosen, type of optimizer) necessary to understand the results?
    \item[] Answer: \answerYes{}.
    \item[] Justification: Section~\ref{sec:setup} reports the model bank, hyperparameters, seeds, splits, normalization, and all three shift protocols; appendix tables give additional protocol variants.

\item {\bf Experiment statistical significance}
    \item[] Question: Does the paper report error bars suitably and correctly defined or other appropriate information about the statistical significance of the experiments?
    \item[] Answer: \answerYes{}.
    \item[] Justification: Main tables and figures report mean $\pm$ population standard deviation over three seeds unless otherwise stated; appendix tables specify when additional mask seeds or target subsets are included.

\item {\bf Experiments compute resources}
    \item[] Question: For each experiment, does the paper provide sufficient information on the computer resources (type of compute workers, memory, time of execution) needed to reproduce the experiments?
    \item[] Answer: \answerYes{}.
    \item[] Justification: Appendix~\ref{sec:appendix_reproducibility} states that runs use one CUDA GPU per job and reports representative logged evaluation times. The paper also gives the frozen-checkpoint protocol so reviewers can verify results without retraining all models.

\item {\bf Code of ethics}
    \item[] Question: Does the research conducted in the paper conform, in every respect, with the NeurIPS Code of Ethics \url{https://neurips.cc/public/EthicsGuidelines}?
    \item[] Answer: \answerYes{}.
    \item[] Justification: The work uses public traffic benchmark data, reports negative and refuted results, and does not involve human subjects, private personal data collection, or deployment decisions.

\item {\bf Broader impacts}
    \item[] Question: Does the paper discuss both potential positive societal impacts and negative societal impacts of the work performed?
    \item[] Answer: \answerYes{}.
    \item[] Justification: Section~\ref{sec:discussion} discusses deployment guidance and limitations. The main positive impact is safer model selection under traffic-sensor failures; the main risk is overgeneralizing the recommendations beyond the evaluated cities and masking model, which the paper explicitly limits.

\item {\bf Safeguards}
    \item[] Question: Does the paper describe safeguards that have been put in place for responsible release of data or models that have a high risk for misuse (e.g., pre-trained language models, image generators, or scraped datasets)?
    \item[] Answer: \answerNA{}.
    \item[] Justification: The paper does not release high-risk generative models, scraped datasets, or sensitive personal data. It uses public traffic forecasting benchmarks and lightweight evaluation artifacts.

\item {\bf Licenses for existing assets}
    \item[] Question: Are the creators or original owners of assets (e.g., code, data, models), used in the paper, properly credited and are the license and terms of use explicitly mentioned and properly respected?
    \item[] Answer: \answerYes{}.
    \item[] Justification: The paper cites the source papers for METR-LA, PEMS-BAY, PEMS04, and all model families. Appendix~\ref{sec:appendix_reproducibility} states that raw datasets are not redistributed and remain under their original terms.

\item {\bf New assets}
    \item[] Question: Are new assets introduced in the paper well documented and is the documentation provided alongside the assets?
    \item[] Answer: \answerYes{}.
    \item[] Justification: The new assets are evaluation scripts, canonical result tables, and a small recommendation artifact rather than a new dataset. Appendix~\ref{sec:appendix_reproducibility} documents their role and file locations.

\item {\bf Crowdsourcing and research with human subjects}
    \item[] Question: For crowdsourcing experiments and research with human subjects, does the paper include the full text of instructions given to participants and screenshots, if applicable, as well as details about compensation (if any)?
    \item[] Answer: \answerNA{}.
    \item[] Justification: The paper does not involve crowdsourcing or human-subject experiments.

\item {\bf Institutional review board (IRB) approvals or equivalent for research with human subjects}
    \item[] Question: Does the paper describe potential risks incurred by study participants, whether such risks were disclosed to the subjects, and whether Institutional Review Board (IRB) approvals (or an equivalent approval/review based on the requirements of your country or institution) were obtained?
    \item[] Answer: \answerNA{}.
    \item[] Justification: The paper does not involve human subjects or participant data collection.

\item {\bf Declaration of LLM usage}
    \item[] Question: Does the paper describe the usage of LLMs if it is an important, original, or non-standard component of the core methods in this research? Note that if the LLM is used only for writing, editing, or formatting purposes and does \emph{not} impact the core methodology, scientific rigor, or originality of the research, declaration is not required.
    \item[] Answer: \answerNA{}.
    \item[] Justification: LLMs are not a component of the scientific method, experiments, routing diagnostic, or recommendation artifact; any language editing assistance does not affect the core methodology or results.

\end{enumerate}
